\scanpage{53}
\mspage{52}

\uline{13 juin} On n'a pas encore exprimé tous les Axiomes \uline{\unsure{Cont|Contr}}\unsure{.} 1, \unsure{....|...}, il \unsure{reste} \unsure{pour compte|par exemple}

\uline{Cont} 3 \unsure{dim} finie

\uline{Cont} 4 ``admissibilité'' -- mais résulte déjà je crois des axiomes C.

\uline{Cont} 6 \inserted{Critère de} disjonction \unsure{via les lieux} -- on y reviendra.

\uline{Cont} 7 Deux lieux de \unsure{str.} dominantes différentes sont disjoints.

Mais avant de continuer, je voudrais encore une fois changer d'optique !

Prendre comme \unsure{ens} de base $\unsure{\mathcal F|\mathscr F}$, l'ens des figures, avec les deux relations d'ordre \[ \leq\ \text{(ou l'inclusion)} \qquad \text{et}\qquad \unsure{\preceq|\preccurlyeq}\ \text{(la subdivision)} \]

\uline{Contrée} (quatrième mouture !)

\unsure{$U$ est|C'est} un ens. $\mathcal F$ d'objets, appelés \uline{figures}, et deux relations d'ordre $\leq$ et $\preceq$ sur $\mathcal F$ : \[ \unsure{\mathcal C|\mathscr C|\mathcal E}=(\unsure{\underline{\mathcal F}|\mathcal F},\leq,\preceq) \] satisfaisant à des axiomes\unsure{,|{}} qu'on va énumérer. Si $F$ et $G$ sont deux figures et $F\leq G$, on dit que $F$ est une \uuline{sous}-\uline{figure}, et on \unsure{écrira|écrit} aussi souvent $F\subset G$. Si $F\preceq G$, on dit que $F$ est une \uline{subdivision} de $G$, ou que $F$ subdivise $G$.
